
When three data attribution methods---TRAK, TracIn, and Kronfluence---are applied to identify which training examples influenced a model's prediction, they can produce substantially different rankings. This raises a question: are these methods measuring the same phenomenon, or do they capture fundamentally different aspects of training data influence?

This question has practical implications. Data attribution underpins applications including identifying training examples that cause model failures, auditing models for fairness violations, and valuing data contributions. If different methods systematically emphasize different aspects of influence---memorization versus feature transfer versus spurious correlation---then method choice shapes the conclusions practitioners draw.

The DATE-LM benchmark \citep{jiao2025datelm} confirmed that no single attribution method dominates across all tasks. TRAK performs well on some evaluations while influence functions perform better on others. However, this finding describes that methods differ without explaining why they differ or what each method measures.

This work investigates whether unexplained disagreement reflects structure: different attribution algorithms compute influence via mathematically distinct operations that create characteristic ''fingerprints.'' TRAK uses random projection of gradients, emphasizing direction alignment in parameter space. TracIn sums checkpoint-weighted gradient dot products, capturing temporal patterns across training. Kronfluence inverts Fisher information via K-FAC approximation, emphasizing curvature in the loss landscape. These operations are not minor variations---they encode different assumptions about what makes training data influential.

A framework is developed for characterizing attribution methods by their mode profiles: systematic patterns of sensitivity to memorization, feature transfer, and spurious association. If methods embed different biases, mode profiles should dissociate---methods should cluster by identity rather than by random variation. If fingerprints are intrinsic to algorithms rather than artifacts of specific datasets, they should exhibit stability within methods.

The contributions are:

First, contrastive mode probing is introduced as a methodology for measuring attribution method sensitivity to specific influence modes using test cases that isolate memorization, feature transfer, and spurious association signals.

Second, quantitative characterization of method dissociation is provided. Using ANOVA analysis across 10 independent runs per method, inter-method variance exceeds intra-method variance with F=1423.55 (threshold: 4.0), and effect sizes are large (Cohen's d=20.64).

Third, transferability of mode profiles across architectures is analyzed. Fingerprints transfer within architectural families (ResNet-ViT r=0.80) but invert across fundamentally different architectures (ConvNeXt shows r=-0.71 to -0.99 versus CNN/Transformer). This finding bounds the generalization of fingerprinting.
